\documentclass[11pt,a4paper]{article}
\usepackage{gtu-el-yazisi}
\begin{document}
\notbilgi{Sayısal Haberleşme}{Sayısal haberleşme sistemi, kanallar ve avantajları}{Ders1.pdf}
\textit{Not: Orijinal dosyadaki basılı sayfalar (s. 2) dijitalleştirilmedi; yalnızca el yazısı sayfalar aktarıldı.}

\sayfa{1}
\section{Sayısal Ders 1}
\noindent Haberleşme: Bir mesajı bir vericiden bir veya birden fazla alıcıya ulaştıran işlem.

\noindent Sayısal işaret: Ayrıkzamanlı, Ayrıkgenlikli sinyal

\noindent Sayısal haberleşme: İşareti sayısal bir şekilde işleyip sonra kanaldan iletmeye denir. (Kanaldaki işaret yine Analogtur.)

\elcizimi[0.9]{sekiller/Ders1-s1-1.png}[1]

\sayfa{3}
\elcizimi[0.7]{sekiller/Ders1-s3-1.png}[3]

\noindent genlik seviyelerin aralığı eşittir normalde

\noindent Kanal kodlaması\\
codeword (kod sözcüğü)

\begin{center}
\begin{tabular}{@{}c@{\quad}c@{\quad}c@{}}
\underline{$\ell=2$} & & \underline{$n=3$}\\[3pt]
00 & $\longrightarrow$ & 000\\
01 & $\longrightarrow$ & 010\\
10 & $\longrightarrow$ & 101\\
11 & $\longrightarrow$ & 111
\end{tabular}
\end{center}
\noindent artık bitler ekleniyor

\[
\begin{array}{c@{\qquad}l}
\ell=2 & g_i(t) \qquad\qquad M=4\\[4pt]
00 & g_1(t) = -2A, \quad 0\le t\le T\\
01 & g_2(t) = -A, \quad ''\\
10 & g_3(t) = A, \quad ''\\
11 & g_4(t) = 2A, \quad ''
\end{array}
\]

\noindent Bandpass mod. M -- phase Shift keying (PSK)
\[ s_i(t) = \sqrt{\frac{2E_s}{T_s}}\cos\!\Bigl(2\pi f_c t + \frac{2\pi}{M}(i-1)\Bigr),\quad 0\le t\le T_s \qquad i=1,\dots,2^{\ell} \]

\noindent \underline{$M=4$}\\
\underline{$\ell=2$}
\[
\begin{array}{ll}
00 & \sqrt{\dfrac{2E_s}{T_s}}\cos(2\pi f_c t) \quad , \quad 0\le t\le T_s\\[8pt]
01 & \sqrt{\dfrac{2E_s}{T_s}}\cos\!\left(2\pi f_c t+\dfrac{\pi}{2}\right) \quad , \quad ''\\[8pt]
10 & \sqrt{\dfrac{2E_s}{T_s}}\cos\!\left(2\pi f_c t+\pi\right) \quad , \quad ''\\[8pt]
11 & \sqrt{\dfrac{2E_s}{T_s}}\cos\!\left(2\pi f_c t+\dfrac{3\pi}{2}\right) \quad , \quad ''
\end{array}
\]

\sayfa{4}
\section{Kanallar}
\subsection{Wireline (kablolu)}
\begin{itemize}
\item Burulu çift (twisted pair) $\sim$kHz
\item Eş eksenli (coaxial) \quad $\sim$MHz
\item Fiber cable \quad $10^{4}$--$10^{5}$ GHz
\item Dalga kılavuzu (Waveguide) \quad 1--100 GHz
\end{itemize}
\subsection{Wireless (Telsiz)}
RF, optik, sualtı kanalında akustik dalga

\noindent Alınan işaret:
\[ r(t) = h(t)\circledast s_i(t) + n(t) \qquad\longrightarrow\ \text{Beyaz Gürültü (AWGN)} \]
\[ \underset{\text{\scriptsize\textcolor{turuncukalem}{kanal idealse}}}{h(t)=\delta(t)} \ \Rightarrow\ r(t) = s_i(t)+n(t) \]
\[ \overset{F}{\longrightarrow}\ R(f) = H(f)\,S_i(f) + N(f) \]

\noindent \textcolor{turuncukalem}{\underline{Bozucu etkiler}}
\begin{itemize}
\item Gürültü (noise)
\item Sönümleme (fading) (çarpımsal gürültü)
\item kanal ideal değilse ($h(t)\neq\delta(t)$) $\rightarrow$ \textcolor{turuncukalem}{Genlik ve faz bozulması (Amplitude and phase distortion)}
\item Karışım (interference)
\item Zayıflama (Attenuation)
\end{itemize}
\[ \textcolor{turuncukalem}{\rightarrow}\ r(t) = \underset{\textcolor{turuncukalem}{\text{\scriptsize sönümleme}}}{a(t)}\,s_i(t) + n(t) \]
\noindent \textcolor{turuncukalem}{Telsiz haberleşmede buna çözüm \okunamadi{üretmek} gerekir}

\elcizimi[0.6]{sekiller/Ders1-s4-1.png}[4]

\noindent Equalization (kanal denkleştirme)

\noindent Kablolu sistemlerde $h(t)$'yi \okunamadi{giderebiliriz}

\noindent \textcolor{kirmizikalem}{$\star$} Mihenk taşı: işaretin Gürültüye Oranı (İGO)\\
Signal to Noise Ratio (SNR)
\formulkutu{\mathrm{SNR} = \frac{P_m}{P_{\tilde n}}}

\noindent \textcolor{kirmizikalem}{$\star$} Bit hata olasılığı (Probability of bit error) (Bit error rate)
\[ \equiv P_b = \mathrm{Prob}(m_i \neq \hat m_i) = f(\mathrm{SNR}) \]
\[ \mathrm{SNR}\uparrow \ \Rightarrow\ P_b\downarrow \qquad \text{(Performans iyileşir)} \]

\sayfa{5}
\noindent \textcolor{kirmizikalem}{$\star$} Kısıtlı kaynaklar: B ve Güç
\section{Sayısal haberleşmenin avantajları}
\noindent 1) Yeniden üretme (Regeneration) \quad \textcolor{kirmizikalem}{$\star$ sayısal haberleşmenin en iyi avantajı}

\elcizimi[0.9]{sekiller/Ders1-s5-1.png}[5]

\noindent 2) Digital circuits are more reliable, cheaper, more flexible

\noindent 3) Better Signal processing (SP) $\longrightarrow$ removing redundancy $\rightarrow$ less BW
\section{Sayısal haberleşmenin dezavantajları}
\noindent 1) Intensive SP as compared to analog systems.

\noindent 2) Synchronization is more complicated as compared to analog systems.

\noindent 3) Nongraceful degradation:\\
When SNR drops below a certain threshold, the quality of service can suddenly change from very good to very poor. In contrast most analog communication systems degrade gracefully.

\end{document}
